\documentclass[12pt,twoside]{report}

%% PREAMBULO
\usepackage[style=apa, backend=biber, sorting=nyt]{biblatex}
\DeclareLanguageMapping{spanish}{spanish-apa}
\addbibresource{referencias.bib}
\usepackage[utf8]{inputenc}
\usepackage[T1]{fontenc}
\usepackage[spanish,es-nodecimaldot]{babel}
\usepackage{lmodern}
\usepackage{amsmath,amssymb,amsthm}
\usepackage{graphicx}
\usepackage[dvipsnames,svgnames,table]{xcolor}
\usepackage[top=2.5cm,bottom=2.5cm,left=3cm,right=2.5cm]{geometry}
\usepackage{fancyhdr}
\usepackage{titlesec}
\usepackage{tikz}
\usepackage{pgfplots}
\pgfplotsset{compat=1.18}
\usetikzlibrary{shapes.geometric,arrows.meta,positioning,fit,calc,
  decorations.pathreplacing,backgrounds,shadows,matrix,chains}
\usepackage{booktabs}
\usepackage{array}
\usepackage{multirow}
\usepackage{tabularx}
\usepackage{longtable}
\usepackage{tcolorbox}
\tcbuselibrary{breakable,skins,theorems}
\usepackage{algorithm}
\usepackage{algpseudocode}
\usepackage{enumitem}
\usepackage{csquotes}
\usepackage{float}
\usepackage{caption}
\usepackage{subcaption}
\usepackage[colorlinks=true,linkcolor=NavyBlue,
            citecolor=ForestGreen,urlcolor=RoyalBlue]{hyperref}

%% COLORES
\definecolor{PrimaryBlue}{RGB}{13,71,161}
\definecolor{AccentBlue}{RGB}{30,136,229}
\definecolor{LightBlue}{RGB}{227,242,253}
\definecolor{DarkGray}{RGB}{40,40,40}
\definecolor{MedGray}{RGB}{100,100,100}
\definecolor{SuccessGreen}{RGB}{46,125,50}
\definecolor{WarnOrange}{RGB}{230,81,0}
\definecolor{PurpleDeep}{RGB}{74,20,140}
\definecolor{TealDark}{RGB}{0,96,100}
\definecolor{BoxBg}{RGB}{245,248,255}

%% SECCIONES
\titleformat{\chapter}[display]
  {\normalfont\huge\bfseries\color{PrimaryBlue}}
  {\chaptertitlename~\thechapter}{10pt}{\Huge}
\titlespacing*{\chapter}{0pt}{-10pt}{20pt}
\titleformat{\section}
  {\normalfont\Large\bfseries\color{PrimaryBlue}}{\thesection}{1em}{}[\titlerule]
\titleformat{\subsection}
  {\normalfont\large\bfseries\color{AccentBlue}}{\thesubsection}{1em}{}
\titleformat{\subsubsection}
  {\normalfont\normalsize\bfseries\color{DarkGray}}{\thesubsubsection}{1em}{}

%% ENCABEZADOS
\pagestyle{fancy}\fancyhf{}
\fancyhead[LE,RO]{\small\color{MedGray}\nouppercase{\leftmark}}
\fancyhead[RE,LO]{\small\color{MedGray}CC3045 -- Inteligencia Artificial}
\fancyfoot[C]{\small\color{MedGray}\thepage}
\renewcommand{\headrulewidth}{0.4pt}

%% CAJAS
\newtcolorbox{definicion}[1][]{enhanced,breakable,
  colback=LightBlue,colframe=PrimaryBlue,
  fonttitle=\bfseries\color{white},title=Definici\'on,
  attach boxed title to top left={yshift=-2mm,xshift=4mm},
  boxed title style={colback=PrimaryBlue,rounded corners},#1}
\newtcolorbox{importante}[1][]{enhanced,breakable,
  colback=BoxBg,colframe=AccentBlue,left=6pt,right=6pt,
  borderline west={3pt}{0pt}{AccentBlue},#1}
\newtcolorbox{modulo}[2][]{enhanced,breakable,
  colback=white,colframe=#1,
  fonttitle=\bfseries,title=#2,
  attach boxed title to top left={yshift=-2mm,xshift=4mm},
  boxed title style={colback=#1,rounded corners}}

%% TEOREMAS
\theoremstyle{definition}
\newtheorem{definition}{Definici\'on}[chapter]
\newtheorem{proposition}{Proposici\'on}[chapter]

%% ALGORITMOS
\renewcommand{\algorithmicrequire}{\textbf{Entrada:}}
\renewcommand{\algorithmicensure}{\textbf{Salida:}}
\algnewcommand\algorithmicforeach{\textbf{para cada}}
\algdef{S}[FOR]{ForEach}[1]{\algorithmicforeach\ #1\ \textbf{hacer}}

\hypersetup{
  pdftitle={Sistema Inteligente H\'ibrido para Planificaci\'on Acad\'emica},
  pdfauthor={Estudiantes UVG},
  pdfsubject={Inteligencia Artificial -- Proyecto Final CC3045}}

%% ============================================================================
\begin{document}

%% PORTADA
\begin{titlepage}
\pagecolor{PrimaryBlue}\color{white}
\begin{tikzpicture}[remember picture,overlay]
  \fill[AccentBlue,opacity=0.4]
    (current page.south west) rectangle
    ([yshift=4cm]current page.south east);
  \fill[AccentBlue,opacity=0.2]
    ([xshift=-2cm,yshift=-2cm]current page.north east) circle(5.5cm);
  \fill[PurpleDeep,opacity=0.15]
    ([xshift=2cm,yshift=-6cm]current page.north east) circle(3.5cm);
  \fill[white,opacity=0.06]
    ([xshift=-4cm,yshift=-4cm]current page.north east) circle(3cm);
\end{tikzpicture}
\vspace*{1.2cm}
\begin{center}
{\fontsize{11}{13}\selectfont\color{white!80}
  UNIVERSIDAD DEL VALLE DE GUATEMALA}\\[0.2cm]
{\fontsize{10}{12}\selectfont\color{white!70}
  Facultad de Ingenier\'ia~--- Ciencias de la Computaci\'on}
\vspace{0.5cm}
\noindent\rule{\linewidth}{0.5pt}
\vspace{0.4cm}

{\fontsize{9}{11}\selectfont\color{white!70}\textsc{Proyecto Final}}\\[0.15cm]
{\fontsize{13}{15}\selectfont\bfseries CC3045~--- Inteligencia Artificial}

\vspace{1.1cm}
{\fontsize{20}{24}\selectfont\bfseries
  Dise\~no de un Sistema Inteligente H\'ibrido\\[0.35cm]
  para la Planificaci\'on Acad\'emica\\[0.35cm]
  Universitaria bajo Restricciones\\[0.35cm]
  e Incertidumbre}

\vspace{0.6cm}
{\fontsize{11}{13}\selectfont\color{white!75}\itshape
  Integraci\'on de aprendizaje supervisado, CSP,\\
  redes bayesianas y MDP para apoyo a la toma\\
  de decisiones acad\'emicas}

\vspace{1.1cm}
\noindent\rule{0.6\linewidth}{0.5pt}
\vspace{0.7cm}

{\fontsize{10}{13}\selectfont
\begin{tabular}{rl}
\textbf{Autores:}  & [Pablo Daniel Barillas Moreno, Hugo Daniel Barillas Ajín y Roberto Nájera]\\[4pt]
\textbf{Carn\'e(s):} & [22193, 23556 y 23781]\\[4pt]
\textbf{Docente:}  & [Ing. Suriano]\\[4pt]
\textbf{Fecha:}    & \today\\[4pt]
\textbf{Semestre:} & I~--- 2026\\
\end{tabular}}
\vspace{1.1cm}

{\fontsize{9}{11}\selectfont\color{white!70}Proyecto Final~}
\end{center}
\end{titlepage}
\nopagecolor

\tableofcontents\clearpage

%% ============================================================================
\chapter*{Resumen}\addcontentsline{toc}{chapter}{Resumen}

\begin{importante}
Este documento presenta el dise\~no, la arquitectura y la metodolog\'ia de
implementaci\'on de un \textbf{sistema inteligente h\'ibrido} para la
planificaci\'on acad\'emica universitaria.  El sistema integra cuatro m\'odulos
de Inteligencia Artificial: (1)~aprendizaje supervisado para la \emph{predicci\'on
de demanda} de cursos; (2)~Problemas de Satisfacci\'on de Restricciones (CSP) con
heur\'isticas para la \emph{generaci\'on de horarios v\'alidos}; (3)~redes
bayesianas para el \emph{razonamiento probabil\'istico} y la estimaci\'on de riesgo;
y (4)~Procesos de Decisi\'on de Markov (MDP) para la \emph{replanificaci\'on
adaptativa} ante contingencias.
\end{importante}

La planificaci\'on acad\'emica en instituciones universitarias es un problema
combinatorio NP-dif\'icil que involucra m\'ultiples capas de decisi\'on, incertidumbre
operativa y restricciones de diversa naturaleza.  Los enfoques tradicionales tratan
estos problemas de manera aislada; este proyecto propone una arquitectura cohesiva
donde la salida de cada m\'odulo alimenta al siguiente, logrando un sistema de
decisi\'on integral.

\textbf{Palabras clave:} CSP, scheduling acad\'emico, redes bayesianas, MDP,
aprendizaje supervisado, heur\'isticas de b\'usqueda, sistema h\'ibrido.

%% ============================================================================
\chapter{Introducci\'on}

\section{Motivaci\'on y contexto}

La planificaci\'on acad\'emica semestral es una de las tareas m\'as cr\'iticas y
complejas en la administraci\'on universitaria.  En una instituci\'on de mediano
tama\~no, el proceso involucra coordinar cientos de cursos, decenas de docentes
con disponibilidades heterog\'eneas, aulas con capacidades y equipamiento diverso,
y restricciones curriculares que deben respetarse para garantizar la progresi\'on
acad\'emica de miles de estudiantes.  Realizarlo manualmente es costoso en tiempo,
propenso a errores y casi siempre suboptimo.

Desde la perspectiva de la Inteligencia Artificial, este problema trasciende la
simple generaci\'on de horarios.  En realidad se trata de un \textbf{problema de
decisi\'on multicapa}: es necesario \emph{predecir} cu\'antas secciones abrir,
\emph{asignar} recursos bajo restricciones, \emph{estimar} el riesgo de conflictos
futuros y \emph{reaccionar} de forma inteligente cuando los planes originales deben
modificarse.

\begin{importante}
Este proyecto demuestra que las t\'ecnicas de IA vistas en el curso~---~agentes
inteligentes, b\'usqueda heur\'istica, CSP, razonamiento probabil\'istico, aprendizaje
supervisado y MDP~---~no son herramientas aisladas sino componentes complementarios
de un sistema de decisi\'on real e integrado.
\end{importante}

\section{Objetivos}

\subsection*{Objetivo general}
Dise\~nar e implementar un sistema inteligente h\'ibrido para la planificaci\'on
acad\'emica universitaria que integre predicci\'on de demanda, generaci\'on de
horarios bajo restricciones, razonamiento probabil\'istico y replanificaci\'on
adaptativa.

\subsection*{Objetivos espec\'ificos}
\begin{enumerate}[leftmargin=2em,itemsep=4pt]
  \item Construir un m\'odulo de predicci\'on de demanda usando aprendizaje supervisado
    y comparar al menos tres algoritmos.
  \item Modelar la generaci\'on de horarios como CSP e implementar backtracking con
    forward checking, MRV y LCV.
  \item Dise\~nar una red bayesiana para estimar la probabilidad de conflictos
    acad\'emicos bajo incertidumbre.
  \item Formalizar un MDP simplificado para la replanificaci\'on adaptativa y
    derivar una pol\'itica de acci\'on.
  \item Integrar los cuatro m\'odulos en un pipeline unificado y evaluarlo sobre
    instancias sint\'eticas.
\end{enumerate}

\section{Estructura del documento}

El cap\'itulo~\ref{ch:problema} describe el problema y sus complejidades.
El cap\'itulo~\ref{ch:analisis} lo analiza desde m\'ultiples perspectivas de IA.
El cap\'itulo~\ref{ch:propuesta} presenta la propuesta global.
El cap\'itulo~\ref{ch:solucion} detalla cada m\'odulo.
El cap\'itulo~\ref{ch:herramientas} documenta las herramientas.
El cap\'itulo~\ref{ch:resultados} presenta metodolog\'ia experimental.
El cap\'itulo~\ref{ch:conclusiones} cierra con conclusiones y trabajo futuro.

%% ============================================================================
\chapter{Descripci\'on del Problema}
\label{ch:problema}

\section{El proceso real de planificaci\'on acad\'emica}

La planificaci\'on acad\'emica semestral en una universidad comprende las siguientes
actividades interrelacionadas:

\begin{enumerate}[leftmargin=2em,itemsep=4pt]
  \item \textbf{Estimaci\'on de demanda}: decidir cu\'antas secciones de cada curso
    deben abrirse, basado en proyecciones de inscripci\'on y capacidad docente.
  \item \textbf{Asignaci\'on de recursos}: asignar a cada secci\'on una franja
    horaria, un aula y un docente, respetando todas las restricciones administrativas
    y acad\'emicas.
  \item \textbf{Gesti\'on de contingencias}: gestionar cambios de \'ultimo momento
    (ausencias, sobrecupo, aver\'ias) de forma r\'apida y coherente.
  \item \textbf{Comunicaci\'on y publicaci\'on}: distribuir el horario final a
    estudiantes, docentes y administraci\'on.
\end{enumerate}

\section{Entradas y salidas del sistema}

\begin{table}[H]
\centering
\caption{Entradas y salidas del sistema de planificaci\'on acad\'emica.}
\label{tab:io}
\begin{tabularx}{\linewidth}{llX}
\toprule
\textbf{Tipo} & \textbf{Entidad} & \textbf{Descripci\'on} \\
\midrule
Entrada & Cursos          & Nombre, cr\'editos, prerrequisitos, carrera \\
Entrada & Estudiantes     & N\'umero proyectado, historial de inscripci\'on \\
Entrada & Docentes        & Disponibilidades, cursos habilitados, carga m\'axima \\
Entrada & Aulas           & Capacidad, equipamiento, disponibilidad \\
Entrada & Hist\'orico     & Inscripciones pasadas, tasas de aprobaci\'on, demanda \\
Entrada & Restricciones   & Reglas duras y blandas administrativas \\
\midrule
Salida  & Predicci\'on    & N\'umero esperado de secciones por curso \\
Salida  & Horario         & Asignaci\'on completa secci\'on$\to$(franja, aula, docente) \\
Salida  & Evaluaci\'on    & M\'etricas de calidad, riesgo de conflicto estimado \\
Salida  & Replanificaci\'on & Acciones sugeridas ante contingencias \\
\bottomrule
\end{tabularx}
\end{table}

\vspace{4cm}

\section{Restricciones del sistema}

\subsection{Restricciones duras}
\begin{itemize}[leftmargin=2em,itemsep=3pt]
  \item Un docente no puede impartir dos secciones en la misma franja horaria.
  \item Un aula no puede albergar dos secciones simult\'aneamente.
  \item La capacidad del aula debe ser mayor o igual al cupo de la secci\'on.
  \item Un docente solo puede impartir cursos para los que est\'a habilitado.
  \item Un docente solo puede asignarse en franjas de su disponibilidad declarada.
\end{itemize}

\subsection{Restricciones blandas}
\begin{itemize}[leftmargin=2em,itemsep=3pt]
  \item Minimizar los intervalos libres en el horario de cada docente.
  \item Respetar las preferencias de franja horaria de los docentes.
  \item Distribuir la carga docente equitativamente entre los d\'ias de la semana.
  \item Concentrar los cursos de un mismo pensum en bloques contiguos.
  \item Evitar asignar cursos de alta demanda en las primeras o \'ultimas franjas.
\end{itemize}

\section{Complejidad del problema}

Sea $n$ el n\'umero de secciones, $t$ el n\'umero de franjas, $r$ el n\'umero de
aulas y $p$ el n\'umero de docentes disponibles.  El espacio de b\'usqueda ingenuo
tiene cardinalidad:

\begin{equation}
  |\mathcal{S}| = (t \cdot r \cdot p)^{n}
  \label{eq:espacio}
\end{equation}

Para una facultad con $n=80$, $t=25$, $r=15$, $p=20$:
$|\mathcal{S}| = 7500^{80} \approx 10^{310}$.
Esto demuestra que la b\'usqueda ingenua es completamente inviable y justifica la
necesidad de m\'etodos inteligentes.

%% ============================================================================
\chapter{An\'alisis del Problema}
\label{ch:analisis}

\section{Descomposici\'on en subproblemas}

El problema global se descompone en cuatro subproblemas con distintas
caracter\'isticas computacionales:

\begin{table}[H]
\centering
\caption{Descomposici\'on del problema y t\'ecnica de IA aplicada.}
\label{tab:descomp}
\begin{tabular}{clll}
\toprule
\textbf{Capa} & \textbf{Subproblema} & \textbf{T\'ecnica de IA} & \textbf{Paradigma} \\
\midrule
1 & Predicci\'on de demanda   & Aprendizaje supervisado & ML \\
2 & Generaci\'on de horarios  & CSP + heur\'isticas     & B\'usqueda \\
3 & Estimaci\'on de riesgo    & Redes bayesianas         & Razonamiento probabil. \\
4 & Replanificaci\'on         & MDP                      & Decisi\'on secuencial \\
\bottomrule
\end{tabular}
\end{table}

\section{El sistema como agente inteligente}

Siguiendo el paradigma de agentes inteligentes de Russell y Norvig~\cite{russell2020},
el sistema puede modelarse como:

\begin{definicion}
Un agente de planificaci\'on acad\'emica es una funci\'on
$f: \mathcal{P}^* \to \mathcal{A}$ que mapea secuencias de percepciones a acciones,
donde:
\begin{itemize}
  \item $\mathcal{P}$ = percepciones: inscripciones hist\'oricas, disponibilidades,
    contingencias.
  \item $\mathcal{A}$ = acciones: asignar secci\'on, mover franja, cancelar secci\'on,
    reasignar docente.
  \item El objetivo es maximizar la m\'etrica de calidad global:
    $Q = \eta_H \cdot \eta_S \cdot (1 - P_{\text{riesgo}})$.
\end{itemize}
\end{definicion}

\begin{figure}[H]
\centering
\begin{tikzpicture}[
  box/.style={rectangle,rounded corners=5pt,draw=PrimaryBlue,thick,
              fill=LightBlue,text width=3.2cm,align=center,minimum height=1cm},
  cloud/.style={ellipse,draw=AccentBlue,thick,fill=blue!10,
                text width=3cm,align=center,minimum height=1.2cm},
  arr/.style={-Stealth,thick,AccentBlue}
]
  \node[cloud] (env) at (0,0) {Entorno\\Acad\'emico};
  \node[box] (sensors) at (5,1.2) {Sensores\\(datos hist\'oricos,\\contingencias)};
  \node[box] (agent) at (5,-1.2) {Agente\\Inteligente\\H\'ibrido};
  \node[box] (actuators) at (10,0) {Actuadores\\(horario,\\replan.)};

  \draw[arr] (env) to[bend left=15] node[above,font=\small]{percepciones} (sensors);
  \draw[arr] (sensors) -- (agent);
  \draw[arr] (agent) -- (actuators);
  \draw[arr] (actuators) to[bend left=15] node[below,font=\small]{acciones} (env);
  \draw[arr,dashed] (agent) to[bend right=20] node[left,font=\small\itshape]{$f(s)$} (sensors);
\end{tikzpicture}
\caption{Modelo del sistema como agente inteligente en su entorno acad\'emico.}
\label{fig:agente}
\end{figure}

\section{An\'alisis de cada subproblema}

\subsection{Capa 1: Predicci\'on como problema de aprendizaje}

La predicci\'on de demanda es una tarea de \textbf{regresi\'on o clasificaci\'on}:
dado un vector de caracter\'isticas hist\'oricas $\mathbf{x}_i$ de un curso,
predecir la demanda esperada $\hat{y}_i$.  Sea $\mathcal{D} = \{(\mathbf{x}_i, y_i)\}_{i=1}^{N}$
el conjunto de entrenamiento.  El objetivo es encontrar una funci\'on $h: \mathbb{R}^d \to \mathbb{R}$
que minimice la funci\'on de p\'erdida:

\begin{equation}
  h^* = \arg\min_{h \in \mathcal{H}} \frac{1}{N}\sum_{i=1}^{N} \mathcal{L}(h(\mathbf{x}_i), y_i) + \lambda \Omega(h)
  \label{eq:erm}
\end{equation}

donde $\mathcal{L}$ es la funci\'on de p\'erdida (MSE para regresi\'on,
entrop\'ia cruzada para clasificaci\'on) y $\Omega(h)$ es un t\'ermino de
regularizaci\'on.

\subsection{Capa 2: Asignaci\'on como CSP}

\begin{definition}
Un CSP se define como $\langle X, D, C \rangle$ donde:
$X = \{X_1,\ldots,X_n\}$ son las variables (secciones),
$D = \{D_1,\ldots,D_n\}$ son los dominios,
y $C$ es el conjunto de restricciones.
\end{definition}

El espacio de estados tiene cardinalidad dada por la ecuaci\'on~\eqref{eq:espacio}.
Las heur\'isticas MRV y LCV, junto con forward checking, reducen este espacio
en varios \'ordenes de magnitud.

\subsection{Capa 3: Incertidumbre como inferencia probabil\'istica}

Los eventos inciertos (ausencia de docente, sobrecupo) se modelan mediante una
red bayesiana con variables latentes.  El teorema de Bayes:

\begin{equation}
  P(Q \mid E=e) = \frac{P(E=e \mid Q)\, P(Q)}{P(E=e)}
  \label{eq:bayes}
\end{equation}

permite actualizar las probabilidades de riesgo dado nuevo evidencia.

\subsection{Capa 4: Reacci\'on como decisi\'on secuencial (MDP)}

El proceso de replanificaci\'on se modela como un MDP
$\langle S, A, T, R, \gamma \rangle$ donde los estados representan
configuraciones del horario y las acciones corresponden a modificaciones
posibles.  La ecuaci\'on de Bellman:

\begin{equation}
  V^*(s) = \max_{a \in A} \left[ R(s,a) + \gamma \sum_{s' \in S} T(s,a,s')\, V^*(s') \right]
  \label{eq:bellman}
\end{equation}

permite derivar la pol\'itica \'optima $\pi^*(s) = \arg\max_{a} Q^*(s,a)$.

%% ============================================================================
\chapter{Propuesta de Soluci\'on}
\label{ch:propuesta}

\section{Arquitectura global del sistema h\'ibrido}

\begin{figure}[H]
\centering
\begin{tikzpicture}[
  mod/.style={rectangle,rounded corners=5pt,draw=#1,thick,fill=#1!10,
              text width=3.4cm,align=center,minimum height=1.3cm,font=\small},
  arrow/.style={-Stealth,thick,gray!70},
  label/.style={font=\scriptsize\itshape,text=MedGray}
]
  %% Datos
  \node[mod=TealDark] (data) at (0,0)
    {\textbf{Datos}\\\footnotesize Hist\'orico, cursos,\\ docentes, aulas};

  %% Modulo 1
  \node[mod=SuccessGreen] (ml) at (4.5,2)
    {\textbf{M\'odulo 1}\\Predicci\'on\\ de Demanda\\(ML)};

  %% Modulo 2
  \node[mod=PrimaryBlue] (csp) at (4.5,0)
    {\textbf{M\'odulo 2}\\Generaci\'on\\ de Horario\\(CSP)};

  %% Modulo 3
  \node[mod=PurpleDeep] (bayes) at (4.5,-2)
    {\textbf{M\'odulo 3}\\Evaluaci\'on\\ de Riesgo\\(Bayes)};

  %% Modulo 4
  \node[mod=WarnOrange] (mdp) at (9.5,0)
    {\textbf{M\'odulo 4}\\Replanifi-\\caci\'on\\(MDP)};

  %% Salida
  \node[mod=TealDark] (out) at (14,0)
    {\textbf{Salida}\\Horario final\\ + Reporte};

  %% Flechas
  \draw[arrow] (data) -- (ml);
  \draw[arrow] (data) -- (csp);
  \draw[arrow] (data) -- (bayes);
  \draw[arrow] (ml) -- node[label,above]{secciones} (csp);
  \draw[arrow] (csp) -- node[label,right]{horario base} (bayes);
  \draw[arrow] (bayes) -- node[label,above]{riesgo} (mdp);
  \draw[arrow] (csp) -- node[label,above]{} (mdp);
  \draw[arrow] (mdp) -- (out);
  \draw[arrow,dashed] (mdp.north) to[bend left=30]
    node[label,above]{retroalimentaci\'on} (csp.north);
\end{tikzpicture}
\caption{Arquitectura global del sistema inteligente h\'ibrido.}
\label{fig:arch}
\end{figure}

\section{Justificaci\'on del enfoque h\'ibrido}

\begin{table}[H]
\centering
\caption{Comparaci\'on entre enfoque \'unico y sistema h\'ibrido propuesto.}
\label{tab:justif}
\begin{tabularx}{\linewidth}{lXX}
\toprule
\textbf{Criterio} & \textbf{Enfoque \'unico (solo CSP)} &
  \textbf{Sistema h\'ibrido propuesto} \\
\midrule
Predicci\'on & No considera demanda futura &
  ML predice secciones necesarias \\
Asignaci\'on & Genera horario bajo restricciones &
  CSP genera horario \'optimo con dominio reducido por ML \\
Incertidumbre & Supone datos perfectos &
  Bayes modela eventos inciertos y estima riesgo \\
Adaptabilidad & Est\'atico ante cambios &
  MDP decide acci\'on \'optima en contingencias \\
Cobertura del curso & Una sola t\'ecnica &
  ML + CSP + Bayes + MDP (4 m\'odulos) \\
Aplicabilidad real & Limitada & Alta \\
\bottomrule
\end{tabularx}
\end{table}

\section{Flujo de decisi\'on del sistema}

El sistema opera secuencialmente en cuatro fases:

\begin{enumerate}[leftmargin=2em,itemsep=5pt]
  \item \textbf{Fase~1 --- Predicci\'on}: el m\'odulo ML analiza el hist\'orico y
    genera una predicci\'on de demanda por curso, determinando cu\'antas secciones
    deben abrirse.
  \item \textbf{Fase~2 --- Generaci\'on}: con las secciones definidas, el m\'odulo CSP
    construye el horario v\'alido aplicando backtracking con MRV, LCV y forward
    checking.
  \item \textbf{Fase~3 --- Evaluaci\'on de riesgo}: la red bayesiana analiza el
    horario generado y estima la probabilidad de cada tipo de conflicto
    ($P(\text{conflicto} \mid \text{horario})$).
  \item \textbf{Fase~4 --- Replanificaci\'on}: si el riesgo supera un umbral
    $\theta$, o si ocurre una contingencia real, el m\'odulo MDP selecciona la
    acci\'on \'optima de reajuste.
\end{enumerate}

%% ============================================================================
\chapter{Descripci\'on de la Soluci\'on}
\label{ch:solucion}

\section{M\'odulo 1: Predicci\'on de Demanda (ML)}

\subsection{Formulaci\'on del problema}

La predicci\'on de demanda se formula como un problema de \textbf{regresi\'on} donde
la variable objetivo es $y_i$ = n\'umero de estudiantes esperados en el curso~$i$
para el siguiente semestre.  Como problema de clasificaci\'on, puede binarizarse:
$y_i = 1$ si se proyecta sobrecupo, 0 en caso contrario.

\subsection{Ingenier\'ia de features}

\begin{table}[H]
\centering
\caption{Variables (features) para el m\'odulo de predicci\'on.}
\label{tab:features}
\begin{tabular}{lll}
\toprule
\textbf{Feature} & \textbf{Tipo} & \textbf{Descripci\'on} \\
\midrule
$x_1$: demanda\_hist & num\'erica & Inscripciones promedio \'ultimos 3 semestres \\
$x_2$: tasa\_aprobacion & num\'erica & Fracci\'on de aprobados histt\'oricamente \\
$x_3$: semestre & categ\'orica & Primer o segundo semestre del a\~no \\
$x_4$: tiene\_prereq & binaria & Si el curso tiene prerrequisito \\
$x_5$: capacidad\_prev & num\'erica & Cupo asignado en el semestre anterior \\
$x_6$: turno & categ\'orica & Ma\~nana / tarde / noche \\
$x_7$: carrera & categ\'orica & Carrera a la que pertenece el curso \\
$x_8$: n\_secciones\_prev & num\'erica & N\'umero de secciones abiertas anteriormente \\
$x_9$: tendencia & num\'erica & Pendiente de la demanda en los \'ultimos 4 semestres \\
\bottomrule
\end{tabular}
\end{table}

\subsection{Modelos comparados}

Se comparan los siguientes modelos de aprendizaje supervisado:
\begin{itemize}[leftmargin=2em,itemsep=3pt]
  \item \textbf{Regresi\'on lineal}: $\hat{y} = \mathbf{w}^\top \mathbf{x} + b$
  \item \textbf{Regresi\'on log\'istica} (versi\'on clasificaci\'on):
    $P(y=1|\mathbf{x}) = \sigma(\mathbf{w}^\top \mathbf{x} + b)$
  \item \textbf{\'Arbol de decisi\'on}: particiones recursivas del espacio de features
  \item \textbf{Naive Bayes}: $\hat{y} = \arg\max_k P(y=k)\prod_j P(x_j|y=k)$
  \item \textbf{SVM}: $\hat{y} = \text{sign}(\mathbf{w}^\top\phi(\mathbf{x}) + b)$
\end{itemize}

\subsection{Pseudoc\'odigo del m\'odulo de predicci\'on}

\begin{algorithm}[H]
\caption{Entrenamiento y evaluaci\'on del m\'odulo de predicci\'on}
\label{alg:ml}
\begin{algorithmic}[1]
\Require Conjunto de datos hist\'orico $\mathcal{D}$, lista de modelos $\mathcal{M}$
\Ensure Mejor modelo $h^*$ y predicciones $\hat{Y}$
\State $\mathcal{D}_{\text{train}}, \mathcal{D}_{\text{test}} \leftarrow \text{SplitStratified}(\mathcal{D}, 0.8)$
\State Preprocesar: imputar, normalizar, codificar categ\'oricas
\State Seleccionar features con correlaci\'on o importancia > umbral
\ForEach{$m \in \mathcal{M}$}
  \State $h_m \leftarrow \text{Entrenar}(m, \mathcal{D}_{\text{train}})$ con validaci\'on cruzada 5-fold
  \State $\text{score}_m \leftarrow \text{Evaluar}(h_m, \mathcal{D}_{\text{test}})$
\EndFor
\State $h^* \leftarrow \arg\max_{m} \text{score}_m$
\State $\hat{Y} \leftarrow h^*(\mathbf{X}_{\text{nuevo semestre}})$
\Return $h^*, \hat{Y}$
\end{algorithmic}
\end{algorithm}

\subsection{M\'etricas de evaluaci\'on del m\'odulo ML}

Para regresi\'on: MAE $= \frac{1}{N}\sum|y_i - \hat{y}_i|$,
RMSE $= \sqrt{\frac{1}{N}\sum(y_i-\hat{y}_i)^2}$.
Para clasificaci\'on: accuracy, precision, recall, F1-score y AUC-ROC.
El criterio de selecci\'on de modelo es:

\begin{equation}
  h^* = \arg\min_{m \in \mathcal{M}} \text{RMSE}_{\text{CV}}(m) + \lambda \cdot \text{Complejidad}(m)
  \label{eq:model_sel}
\end{equation}

%% ------------- SECCION 2 CSP -------------------------------------------------
\section{M\'odulo 2: Generaci\'on de Horarios con CSP}

\subsection{Formulaci\'on formal}

\begin{definition}
El CSP de scheduling acad\'emico es $\langle X, D, C \rangle$ donde:
\begin{itemize}
  \item $X_i = (\text{franja}_i, \text{aula}_i, \text{docente}_i)$ para cada secci\'on $i$.
  \item $D_i = T \times R_i \times P_i$ donde $T$ es el conjunto de franjas,
    $R_i$ las aulas con capacidad suficiente, y $P_i$ los docentes habilitados.
  \item $C$ incluye las restricciones de no solapamiento (ecuaciones
    \ref{eq:rc1}--\ref{eq:rc3}).
\end{itemize}
\end{definition}

\begin{equation}
  \forall i \neq j:\; (\text{doc}_i = \text{doc}_j) \Rightarrow (\text{franja}_i \neq \text{franja}_j)
  \label{eq:rc1}
\end{equation}
\begin{equation}
  \forall i \neq j:\; (\text{aula}_i = \text{aula}_j) \Rightarrow (\text{franja}_i \neq \text{franja}_j)
  \label{eq:rc2}
\end{equation}
\begin{equation}
  \forall i:\; \text{cap}(\text{aula}_i) \geq \text{cupo}(i)
  \label{eq:rc3}
\end{equation}

\subsection{Factor graph y funci\'on objetivo}

\begin{equation}
  \text{Weight}(\mathbf{x}) = \prod_{j} f_j(\mathbf{x}_{\text{scope}(j)}), \quad f_j \in [0,1]
  \label{eq:weight}
\end{equation}

\begin{equation}
  \mathbf{x}^* = \arg\max_{\mathbf{x} \in \mathcal{F}} \text{Weight}(\mathbf{x})
  \label{eq:obj_csp}
\end{equation}

donde $\mathcal{F}$ es el conjunto de asignaciones que satisfacen todas las restricciones duras.

\subsection{Grafo de restricciones}

\begin{figure}[H]
\centering
\begin{tikzpicture}[
  var/.style={circle,draw=PrimaryBlue,thick,fill=LightBlue,
              minimum size=1cm,font=\small},
  res/.style={rectangle,draw=WarnOrange,thick,fill=orange!10,
              rounded corners=3pt,font=\scriptsize,align=center,
              minimum height=0.7cm,minimum width=1.8cm}
]
  \node[var] (s1) at (0,0)  {$X_1$};
  \node[var] (s2) at (3,0)  {$X_2$};
  \node[var] (s3) at (6,0)  {$X_3$};
  \node[var] (s4) at (1.5,-2.5) {$X_4$};
  \node[var] (s5) at (4.5,-2.5) {$X_5$};

  \draw[thick,AccentBlue] (s1) -- node[res,above,yshift=2pt]{docente} (s2);
  \draw[thick,AccentBlue] (s2) -- node[res,above,yshift=2pt]{aula}    (s3);
  \draw[thick,AccentBlue] (s1) -- node[res,left, xshift=-3pt]{pensum} (s4);
  \draw[thick,AccentBlue] (s3) -- node[res,right,xshift=3pt]{docente} (s5);
  \draw[thick,AccentBlue] (s4) -- node[res,below,yshift=-2pt]{aula}   (s5);
  \draw[thick,dashed,MedGray] (s2) -- (s4);
  \draw[thick,dashed,MedGray] (s2) -- (s5);
\end{tikzpicture}
\caption{Grafo de restricciones entre secciones.  Aristas s\'olidas = restricciones activas.}
\label{fig:grafo}
\end{figure}

\subsection{Heur\'isticas de b\'usqueda}

\paragraph{MRV (Minimum Remaining Values):}
\begin{equation}
  X^* = \arg\min_{X_i \notin \text{asignadas}} |D_i^{\text{actual}}|
  \label{eq:mrv}
\end{equation}

\paragraph{LCV (Least Constraining Value):}
\begin{equation}
  \text{lcv}(v, X_i) = \sum_{X_j \in N(X_i)} |\{d \in D_j : (v,d)\text{ viola }C\}|
  \label{eq:lcv}
\end{equation}

\begin{algorithm}[H]
\caption{Backtracking con FC, MRV y LCV}
\label{alg:bt}
\begin{algorithmic}[1]
\Require CSP $\langle X,D,C\rangle$
\Ensure Asignaci\'on completa o FALLO
\Function{Backtrack}{$A$, $D$}
  \If{$A$ completa} \Return $A$ \EndIf
  \State $X_i \leftarrow \text{MRV}(A, D)$
  \ForEach{$v \in \text{LCV}(X_i, D, A)$}
    \If{$v$ consistente con $A$}
      \State $A \leftarrow A \cup \{X_i = v\}$
      \State $D' \leftarrow \text{ForwardChecking}(X_i, v, D)$
      \If{$D' \neq$ FALLO}
        \State resultado $\leftarrow$ \Call{Backtrack}{$A$, $D'$}
        \If{resultado $\neq$ FALLO} \Return resultado \EndIf
      \EndIf
      \State $A \leftarrow A \setminus \{X_i=v\}$; restaurar $D$
    \EndIf
  \EndFor
  \Return FALLO
\EndFunction
\end{algorithmic}
\end{algorithm}

\begin{algorithm}[H]
\caption{Beam Search para generaci\'on de horarios}
\label{alg:beam}
\begin{algorithmic}[1]
\Require Ancho $k$, CSP, funci\'on de evaluaci\'on $\text{Eval}$
\Ensure Mejor asignaci\'on encontrada
\Function{BeamSearch}{$k$, CSP}
  \State beam $\leftarrow [\{\}]$
  \While{beam no vac\'io}
    \State candidatos $\leftarrow []$
    \ForEach{$A \in$ beam}
      \State $X_i \leftarrow$ siguiente variable no asignada (MRV)
      \ForEach{$v \in D_i$}
        \If{$A \cup \{X_i=v\}$ consistente}
          candidatos.add$(A \cup \{X_i=v\})$
        \EndIf
      \EndFor
    \EndFor
    \If{candidatos vac\'io} \Return mejor en beam \EndIf
    \State beam $\leftarrow$ top-$k$ de candidatos seg\'un Eval
  \EndWhile
\EndFunction
\end{algorithmic}
\end{algorithm}

%% ------------- SECCION 3 BAYES -----------------------------------------------
\section{M\'odulo 3: Evaluaci\'on de Riesgo con Redes Bayesianas}

\subsection{Motivaci\'on}

Una vez generado el horario base, existen eventos inciertos que podr\'ian
invalidarlo: ausencia inesperada de un docente, demanda mayor a la predicha,
falla de equipo en un aula especializada.  La red bayesiana permite cuantificar
\emph{a priori} la probabilidad de cada tipo de conflicto.

\subsection{Estructura de la red bayesiana}

\begin{figure}[H]
\centering
\begin{tikzpicture}[
  bnode/.style={ellipse,draw=PurpleDeep,thick,fill=purple!10,
                minimum width=2.4cm,minimum height=0.9cm,
                align=center,font=\small},
  barr/.style={-Stealth,thick,PurpleDeep}
]
  \node[bnode] (demand) at (0, 3)    {Alta\\demanda};
  \node[bnode] (tdoc)   at (4, 3)    {Disponib.\\docente};
  \node[bnode] (taula)  at (8, 3)    {Disponib.\\aula};
  \node[bnode] (sobre)  at (2, 1.2)  {Riesgo de\\sobrecupo};
  \node[bnode] (change) at (6, 1.2)  {Riesgo de\\cambio};
  \node[bnode] (reprog) at (4,-0.4)  {Riesgo de\\replanif.};

  \draw[barr] (demand) -- (sobre);
  \draw[barr] (demand) -- (change);
  \draw[barr] (tdoc)   -- (change);
  \draw[barr] (taula)  -- (change);
  \draw[barr] (sobre)  -- (reprog);
  \draw[barr] (change) -- (reprog);
\end{tikzpicture}
\caption{Red bayesiana para estimaci\'on de riesgo acad\'emico.}
\label{fig:bayes}
\end{figure}

\subsection{Inferencia probabil\'istica}

La probabilidad de que se requiera replanificaci\'on dada la evidencia $E=e$:

\begin{equation}
  P(\text{Replanif.} = 1 \mid E = e) =
  \frac{P(E=e \mid \text{Replanif.}=1)\; P(\text{Replanif.}=1)}{P(E=e)}
  \label{eq:bayes_risk}
\end{equation}

El algoritmo de inferencia exacta (eliminaci\'on de variables) calcula
la distribuci\'on marginal de cualquier nodo dado un conjunto de nodos
observados.

\begin{algorithm}[H]
\caption{Evaluaci\'on de riesgo probabil\'istico}
\label{alg:risk}
\begin{algorithmic}[1]
\Require Red bayesiana $\mathcal{B}$, horario $H$, umbral $\theta$
\Ensure Probabilidad de riesgo $p_{\text{risk}}$ y bandera de alerta
\State Extraer evidencia $e$ del horario $H$: disponibilidades, demanda predicha
\State $p_{\text{risk}} \leftarrow P(\text{Replanif.}=1 \mid e)$ (inferencia en $\mathcal{B}$)
\If{$p_{\text{risk}} > \theta$}
  \Return $p_{\text{risk}}$, ALERTA
\Else
  \Return $p_{\text{risk}}$, OK
\EndIf
\end{algorithmic}
\end{algorithm}

%% ------------- SECCION 4 MDP ------------------------------------------------
\section{M\'odulo 4: Replanificaci\'on con MDP}

\subsection{Formulaci\'on del MDP}

Cuando ocurre una contingencia, el proceso de replanificaci\'on se formaliza
como un MDP $\langle S, A, T, R, \gamma \rangle$:

\begin{table}[H]
\centering
\caption{Componentes del MDP de replanificaci\'on.}
\label{tab:mdp}
\begin{tabular}{lp{10cm}}
\toprule
\textbf{Componente} & \textbf{Descripci\'on} \\
\midrule
$S$ & Configuraciones del horario: (secciones asignadas, conflictos activos,
  recursos disponibles) \\
$A$ & Acciones: mover franja, reasignar aula, cambiar docente, abrir nueva
  secci\'on, mantener plan \\
$T(s,a,s')$ & Probabilidad de transitar al estado $s'$ al tomar acci\'on $a$
  en estado $s$ (derivada de la red bayesiana) \\
$R(s,a)$ & Recompensa: $+10$ si se resuelve el conflicto sin afectar otros
  cursos; $-5$ si genera nuevos conflictos \\
$\gamma$ & Factor de descuento $\gamma = 0.9$ \\
\bottomrule
\end{tabular}
\end{table}

\subsection{Ecuaci\'on de Bellman}

La funci\'on de valor \'optima satisface:

\begin{equation}
  V^*(s) = \max_{a \in A(s)} \left[ R(s,a) + \gamma \sum_{s' \in S} T(s,a,s')\, V^*(s') \right]
  \label{eq:bellman2}
\end{equation}

La pol\'itica \'optima se extrae como:

\begin{equation}
  \pi^*(s) = \arg\max_{a \in A(s)} \left[ R(s,a) + \gamma \sum_{s'} T(s,a,s')\, V^*(s') \right]
  \label{eq:policy}
\end{equation}

\subsection{Diagrama de transici\'on de estados}

\begin{figure}[H]
\centering
\begin{tikzpicture}[
  state/.style={circle,draw=WarnOrange,thick,fill=orange!15,
                minimum size=1.5cm,align=center,font=\small},
  sarr/.style={-Stealth,thick,WarnOrange},
  lbl/.style={font=\scriptsize,color=DarkGray}
]
  \node[state] (s0) at (0,0) {$s_0$\\Normal};
  \node[state] (s1) at (4,1.5) {$s_1$\\Conflicto\\ leve};
  \node[state] (s2) at (4,-1.5) {$s_2$\\Ausencia\\ docente};
  \node[state] (s3) at (8,0) {$s_3$\\Resuelto};
  \node[state,draw=SuccessGreen,fill=green!15] (sf) at (12,0)
    {$s_f$\\Horario\\publicado};

  \draw[sarr] (s0) to[bend left=10]
    node[lbl,above]{$P(\text{leve})$} (s1);
  \draw[sarr] (s0) to[bend right=10]
    node[lbl,below]{$P(\text{ausencia})$} (s2);
  \draw[sarr] (s1) to[bend left=10]
    node[lbl,above]{$a$: mover franja} (s3);
  \draw[sarr] (s2) to[bend right=10]
    node[lbl,below]{$a$: reasignar doc.} (s3);
  \draw[sarr] (s3) -- node[lbl,above]{verificar} (sf);
  \draw[sarr,dashed] (s1) to[bend right=15]
    node[lbl,right]{$a$: abrir secc.} (s2);
\end{tikzpicture}
\caption{Diagrama simplificado de transici\'on de estados en el MDP de replanificaci\'on.}
\label{fig:mdp}
\end{figure}

\begin{algorithm}[H]
\caption{Selecci\'on de acci\'on en el MDP (pol\'itica greedy)}
\label{alg:mdp}
\begin{algorithmic}[1]
\Require Estado actual $s$, funci\'on $Q^*(s,a)$, MDP
\Ensure Acci\'on \'optima $a^*$
\State Observar estado del horario: conflictos, recursos, probabilidades de riesgo
\State Actualizar $s \in S$ con el estado actual observado
\State $a^* \leftarrow \arg\max_{a \in A(s)} Q^*(s,a)$
\State Ejecutar $a^*$ sobre el horario
\State Observar nuevo estado $s'$ y recompensa $r$
\State Actualizar estimaciones si se usa aprendizaje por refuerzo
\Return $a^*$, $s'$
\end{algorithmic}
\end{algorithm}

%% ------------- SECCION 5 INTEGRACION ----------------------------------------
\section{Integraci\'on de m\'odulos}

\subsection{Pipeline completo de decisi\'on}

\begin{figure}[H]
\centering
\begin{tikzpicture}[
  step/.style={rectangle,rounded corners=4pt,draw=#1,thick,fill=#1!12,
               text width=3cm,align=center,minimum height=1.1cm,font=\small},
  arr/.style={-Stealth,thick,gray}
]
  \node[step=TealDark] (d) at (0,0) {Datos hist\'oricos\\+ restricciones};
  \node[step=SuccessGreen] (m1) at (3.5,0) {M\'od. 1\\ML: predicci\'on\\de secciones};
  \node[step=PrimaryBlue] (m2) at (7,0) {M\'od. 2\\CSP: horario\\base};
  \node[step=PurpleDeep] (m3) at (10.5,0) {M\'od. 3\\Bayes: riesgo\\$P(\text{conf})$};
  \node[step=WarnOrange] (m4) at (7,-2.5) {M\'od. 4\\MDP: acci\'on\\de reajuste};
  \node[step=TealDark] (out) at (3.5,-2.5) {Horario final\\+ reporte};

  \draw[arr] (d)--(m1)--(m2)--(m3);
  \draw[arr] (m3) -- node[right,font=\scriptsize]{riesgo$>\theta$} (m4);
  \draw[arr] (m4) -- node[below,font=\scriptsize]{horario corregido} (out);
  \draw[arr] (m2) to[bend right=20] node[below,font=\scriptsize]{OK} (out);
  \draw[arr,dashed] (m4.north) to[bend left=40]
    node[above,font=\scriptsize\itshape]{retroalimentaci\'on} (m2.south);
\end{tikzpicture}
\caption{Pipeline completo de decisi\'on del sistema h\'ibrido.}
\label{fig:pipeline}
\end{figure}

\subsection{Funci\'on objetivo global}

La calidad global del sistema se define como:

\begin{equation}
  Q_{\text{global}} = \underbrace{\eta_H}_{\text{CSP duras}} \cdot
  \underbrace{\eta_S}_{\text{CSP blandas}} \cdot
  \underbrace{(1 - P_{\text{riesgo}})}_{\text{Bayes}} \cdot
  \underbrace{\frac{V^*(s_{\text{init}})}{V_{\max}}}_{\text{MDP}}
  \label{eq:global}
\end{equation}

El sistema busca maximizar $Q_{\text{global}}$ a trav\'es del pipeline de cuatro m\'odulos.

%% ============================================================================
\chapter{Herramientas Aplicadas}
\label{ch:herramientas}

\section{Stack tecnol\'ogico}

\begin{table}[H]
\centering
\caption{Herramientas y bibliotecas del proyecto.}
\begin{tabular}{llp{7cm}}
\toprule
\textbf{Herramienta} & \textbf{Versi\'on} & \textbf{Uso} \\
\midrule
Python & 3.11 & Lenguaje principal \\
Jupyter Notebook & 7.x & Documentaci\'on ejecutable (.ipynb) \\
NumPy / Pandas & 1.26 / 2.1 & Manejo de datos y matrices \\
scikit-learn & 1.4 & Modelos ML, CV, m\'etricas \\
pgmpy & 0.1.x & Redes bayesianas e inferencia \\
Matplotlib / Seaborn & 3.8 & Visualizaci\'on \\
\texttt{constraint} & 1.4 & Biblioteca CSP de referencia \\
Overleaf & --- & Documentaci\'on LaTeX \\
Git / GitHub & --- & Repositorio p\'ublico y control de versiones \\
\bottomrule
\end{tabular}
\end{table}

\section{Uso de IA Generativa}
\label{sec:ia}

\begin{importante}
Esta secci\'on documenta, conforme a las instrucciones del curso, el uso de
herramientas de IA generativa como apoyo al desarrollo del proyecto.
\end{importante}

\subsection{Prompts utilizados}

\begin{enumerate}[leftmargin=2em,itemsep=8pt]

  \item \textbf{Prompt de exploraci\'on del dise\~no h\'ibrido:}
  \begin{quote}\itshape
  ``Necesito dise\~nar un sistema de planificaci\'on acad\'emica que use m\'as de una
  t\'ecnica de IA.  Qu\'e combinaci\'on de m\'etodos ser\'ia coherente y justificable?
  Quiero que se relacione con: aprendizaje supervisado, CSP, redes bayesianas y MDP.''
  \end{quote}
  \textbf{Por qu\'e funcion\'o:} el prompt especific\'o expl\'icitamente las cuatro
  t\'ecnicas objetivo, lo que orient\'o la respuesta hacia integraci\'on real en lugar
  de descripci\'on aislada de cada m\'etodo.

  \vspace{3cm}

  \item \textbf{Prompt de formalizaci\'on del MDP:}
  \begin{quote}\itshape
  ``Dame un ejemplo concreto de c\'omo modelar la replanificaci\'on de horarios como
  un MDP: define estados, acciones, recompensas y la ecuaci\'on de Bellman aplicada.''
  \end{quote}
  \textbf{Por qu\'e funcion\'o:} solicitar un ejemplo concreto con los componentes
  formales oblig\'o a la IA a producir definiciones precisas, verificables contra la
  bibliograf\'ia del curso.

  \item \textbf{Prompt de depuraci\'on del m\'odulo CSP:}
  \begin{quote}\itshape
  ``Mi forward checking no est\'a propagando correctamente cuando hay m\'as de un docente
  elegible por secci\'on.  Revisa este c\'odigo y explica el error conceptual.''
  \end{quote}
  \textbf{Por qu\'e funcion\'o:} proporcionar el c\'odigo real permiti\'o un
  diagn\'ostico preciso; la IA no invent\'o el problema.

\end{enumerate}

\subsection{Evoluci\'on de la idea}

\begin{description}[leftmargin=2em,itemsep=4pt]
  \item[\textbf{Iteraci\'on 1.}] Inicio con solo CSP. La IA help\'o a identificar
    que el problema era m\'as rico e incluir predicci\'on y riesgo.
  \item[\textbf{Iteraci\'on 2.}] Incorporaci\'on del m\'odulo ML para predicci\'on
    de demanda; la IA sugiri\'o el pipeline de datos y la selecci\'on de features.
  \item[\textbf{Iteraci\'on 3.}] Dise\~no de la red bayesiana; la IA explic\'o la
    estructura de nodos; la topolog\'ia final fue dise\~nada y validada por el equipo.
  \item[\textbf{Iteraci\'on 4.}] Formalizaci\'on del MDP; ecuaciones y pseudoc\'odigo
    fueron desarrollados por el equipo con la IA como referencia.
\end{description}

\subsection{Delimitaci\'on entre asistencia IA e implementaci\'on propia}

\begin{table}[H]
\centering
\begin{tabular}{p{7cm}cc}
\toprule
\textbf{Componente} & \textbf{Asistencia IA} & \textbf{Propio} \\
\midrule
Dise\~no del pipeline h\'ibrido & Orientaci\'on & \checkmark \\
Modelado formal CSP             & Bibliograf\'ia & \checkmark \\
Implementaci\'on backtracking + FC & Revisi\'on & \checkmark \\
Topolog\'ia de la red bayesiana  & Sugerencia & \checkmark \\
Formalizaci\'on MDP y Bellman    & Referencia & \checkmark \\
Features de predicci\'on (ML)    & Lista inicial & \checkmark \\
Dise\~no de experimentos          & Sugerencia & \checkmark \\
Redacci\'on acad\'emica           & Apoyo estilo & \checkmark \\
\bottomrule
\end{tabular}
\end{table}

\section{Repositorio y entregables complementarios}

Todo el c\'odigo fuente est\'a disponible en un repositorio p\'ublico:
\begin{center}
\noindent\textbf{Repositorio del proyecto final:} \href{https://github.com/DanielBarillasM/Proyecto-Final_IA_Seccion_10_Grupo_1}{GitHub}
\end{center}

\begin{enumerate}[leftmargin=2em,itemsep=3pt]
  \item \textbf{Documento PDF} (este archivo).
  \item \textbf{Notebook Jupyter} (\texttt{.ipynb}) con implementaci\'on y experimentos.
  \item \textbf{C\'odigo fuente} organizado por m\'odulos en el repositorio.
  \item \textbf{Video explicativo} ($\approx$7~min) con demostraci\'on del sistema.
  \item \textbf{Presentaci\'on} (m\'ax. 10 diapositivas) para la exposici\'on.
\end{enumerate}

%% ============================================================================
\chapter{Resultados y M\'etricas}
\label{ch:resultados}

\section{Metodolog\'ia experimental}

La evaluaci\'on del sistema sigue un protocolo por m\'odulo.
El experimento utiliz\'o el dataset \textbf{OULAD real}
con $\mathbf{n=11\,999}$ observaciones (datos de la Open University, UK).
Los modelos ML se evaluaron con un 25\,\% de datos de prueba~($n_{\text{test}}=3\,000$).
El CSP resolvi\'o una instancia de \textbf{10 secciones} con 25~franjas, 7~aulas
y 8~profesores.
La red bayesiana y el MDP fueron evaluados sobre ese mismo horario generado.

\begin{enumerate}[leftmargin=2em,itemsep=4pt]
  \item \textbf{M\'odulo ML}: comparaci\'on de 5 modelos (Random Forest, Extra Trees,
    Regresi\'on log\'istica, SGD y baseline); tiempo total de entrenamiento: $0.71$~s.
  \item \textbf{M\'odulo CSP}: instancia Tiny ($n=10$) resuelta con BT+FC+MRV+LCV;
    0 retrocesos, $\eta_H=1.0$.
  \item \textbf{M\'odulo Bayes}: red bayesiana de 6 nodos; evaluada sobre las
    10 secciones asignadas; Brier score~$=0.117$.
  \item \textbf{M\'odulo MDP}: 5 estados, 5 acciones; pol\'itica derivada con
    value~iteration ($\gamma=0.9$); $V^*(s_{\text{normal}})=140.0$.
\end{enumerate}

\begin{importante}
\textbf{KPIs globales del experimento real:}
\begin{itemize}[itemsep=2pt]
  \item Fuente de datos: OULAD real ($n=11\,999$, modo r\'apido).
  \item Mejor modelo: \textbf{Random Forest} --- F1\,=\,0.934, AUC\,=\,0.981.
  \item Umbral de decisi\'on \'optimo: $\tau=0.23$ (m\'inimo coste FP/FN).
  \item CSP: 10/10 secciones asignadas, \textbf{0 retrocesos}, tiempo~$<0.01$~s.
  \item Red bayesiana: riesgo medio $P(\text{Replanif})=0.719$; Brier~$=0.117$.
  \item MDP: $V^*(s_{\text{normal}})=140.0$; acci\'on dominante: \textit{open\_extra}.
  \item $Q_{\text{global}}$ sistema completo: $\mathbf{0.219}$.
\end{itemize}
\end{importante}

\section{M\'etricas por m\'odulo}

\begin{table}[H]
\centering
\caption{M\'etricas de evaluaci\'on por m\'odulo con resultados reales del experimento.}
\label{tab:metricas}
\begin{tabular}{lllll}
\toprule
\textbf{M\'odulo} & \textbf{M\'etrica} & \textbf{Definici\'on} & \textbf{Objetivo} & \textbf{Resultado} \\
\midrule
ML & F1-score      & $2\cdot\frac{P\cdot R}{P+R}$                    & Maximizar & \textbf{0.934} \\
ML & AUC-ROC       & \'Area bajo la curva ROC                         & Maximizar & \textbf{0.981} \\
ML & Accuracy      & Fracci\'on de predicciones correctas             & Maximizar & 0.933 \\
ML & Umbral $\tau$ & Umbral de decisi\'on \'optimo (coste)            & ---       & 0.23 \\
CSP & $\eta_H$     & Fracci\'on restricc. duras satisfechas           & $=1.0$    & \textbf{1.00} \\
CSP & $\eta_S$     & Calidad de restricc. blandas                     & Maximizar & 1.00$^*$ \\
CSP & Backtracks   & N\'um.\ de retrocesos                            & Minimizar & \textbf{0} \\
CSP & $T$ (s)      & Tiempo de c\'omputo                              & Minimizar & ${<}0.01$ \\
Bayes & Log-lik.   & $\sum \log P(e_i|\mathcal{B})$                   & Maximizar & $-27.16$ \\
Bayes & Brier score& $\frac{1}{N}\sum(p_i-o_i)^2$                    & Minimizar & \textbf{0.117} \\
Bayes & $P(\text{Replanif})$ & Riesgo medio de replanificaci\'on       & Referencia & 0.719 \\
MDP & $V^*(s_0)$   & Valor del estado normal bajo $\pi^*$             & Maximizar & \textbf{140.0} \\
MDP & Acci\'on MRV & Acci\'on m\'as frecuente seg\'un pol\'itica      & ---       & \textit{open\_extra} \\
Global & $Q_{\text{global}}$ & Ec.~\eqref{eq:global}                  & Maximizar & 0.219 \\
\bottomrule
\multicolumn{5}{l}{\small $^*$ Inferida: ninguna secci\'on requiri\'o retroceso; cada asignaci\'on
  obtuvo la primera opci\'on de mayor puntaje.}
\end{tabular}
\end{table}

\section{Instancias de prueba}

\begin{table}[H]
\centering
\caption{Instancias de prueba para el m\'odulo CSP.
  La instancia \textit{Tiny} corresponde a la ejecuci\'on real del notebook
  ($n_{\text{secc.}}=10$, $t=25$ franjas, $r=7$ aulas, $p=8$ profesores).
  Las instancias Small--Large son proyecciones no evaluadas en este experimento.}
\label{tab:inst}
\begin{tabular}{ccccccccc}
\toprule
\textbf{Inst.} & $n$ & $t$ & $r$ & $p$ &
  $\eta_H$ & $\eta_S$ & Backtracks & $T$(s) \\
\midrule
Tiny   & 10  & 25 & 7  & 8  & 1.00 & 1.00$^*$ & 0 & ${<}0.01$ \\
Small  & 25  & 25 & 7  & 8  & ---  & ---       & ---        & --- \\
Medium & 50  & 25 & 7  & 8  & ---  & ---       & ---        & --- \\
Large  & 100 & 25 & 7  & 8  & ---  & ---       & ---        & --- \\
\bottomrule
\multicolumn{9}{l}{\small $^*$ $\eta_S$ inferida: sin retrocesos, toda secci\'on
  asignada a su primera opci\'on de mayor puntaje.}
\end{tabular}
\end{table}

\section{Comparaci\'on de modelos ML}

\begin{table}[H]
\centering
\caption{Comparaci\'on de modelos para predicci\'on de riesgo acad\'emico
  (dataset OULAD real, $n=11\,999$ observaciones, conjunto de prueba 25\,\%).
  El mejor modelo por F1 se destaca en negrita.}
\label{tab:ml_comp}
\begin{tabular}{lccccc}
\toprule
\textbf{Modelo} & \textbf{Accuracy} & \textbf{Precisi\'on} & \textbf{Recall} & \textbf{F1} & \textbf{AUC} \\
\midrule
\textbf{Random Forest}    & \textbf{0.933} & \textbf{0.972} & \textbf{0.900} & \textbf{0.934} & \textbf{0.981} \\
Extra Trees               & 0.929 & 0.967 & 0.895 & 0.930 & 0.974 \\
Regresi\'on log\'istica   & 0.871 & 0.896 & 0.855 & 0.875 & 0.942 \\
SGD (log-loss)            & 0.858 & 0.835 & 0.910 & 0.871 & 0.939 \\
Dummy (baseline)          & 0.528 & 0.528 & 1.000 & 0.691 & 0.500 \\
\bottomrule
\end{tabular}
\end{table}

\section{Comparaci\'on de estrategias de b\'usqueda CSP}

\begin{figure}[H]
\centering
\begin{tikzpicture}
\begin{axis}[
  width=12cm,height=6cm,
  xlabel={Tama\~no de instancia ($n$)},
  ylabel={Backtracks (escala log)},
  ymode=log,
  xtick={10,25,50,100},
  xticklabels={Tiny,Small,Medium,Large},
  legend pos=north west,
  grid=major,grid style={dashed,gray!30},thick,
  title={BT+FC+MRV+LCV: resultado real (Tiny, $n\!=\!10$, backtracks$=\!0$);
    curvas Small--Large son proyecciones te\'oricas},
  title style={font=\small\itshape,color=MedGray}
]
\addplot[color=WarnOrange,mark=square*,dashed]
  coordinates{(10,12)(25,380)(50,15000)(100,2100000)};
\addlegendentry{BT puro (proyectado)}
\addplot[color=AccentBlue,mark=*,thick]
  coordinates{(10,1)(25,42)(50,310)(100,4200)};
\addlegendentry{BT+FC+MRV+LCV (\textbf{real}: 0 en $n\!=\!10$, proyectado para $n>10$)}
\addplot[color=SuccessGreen,mark=triangle*,thick]
  coordinates{(10,0.1)(25,0.1)(50,0.1)(100,0.1)};
\addlegendentry{Beam Search ($k=5$, proyectado)}
% Punto real medido
\addplot[color=AccentBlue,mark=*,mark size=4pt,only marks]
  coordinates{(10,1)};
\addlegendentry{Resultado real (backtracks$=\!0$, graficado como 1 por escala log)}
\end{axis}
\end{tikzpicture}
\caption{Comparaci\'on ilustrativa de backtracks por estrategia de b\'usqueda.}
\label{fig:bt_comp}
\end{figure}

\section{Ablation study}

\begin{table}[H]
\centering
\caption{Ablation study del sistema h\'ibrido con valores calculados sobre el experimento real.
  $Q_{\text{global}} = \eta_H \cdot \eta_S \cdot (1-P_{\text{riesgo}}) \cdot V^*(s_0)/V_{\max}$,
  donde los factores omitidos se fijan en~1.
  Datos: $\eta_H\!=\!1.00$, $\eta_S\!=\!1.00$, $P_{\text{riesgo}}\!=\!0.719$,
  $V^*(s_{\text{normal}})\!=\!140.0$, $V_{\max}\!=\!180.0$.}
\label{tab:ablation}
\begin{tabular}{lccccc}
\toprule
\textbf{Configuraci\'on} & \textbf{ML} & \textbf{CSP} & \textbf{Bayes} & \textbf{MDP} &
  $Q_{\text{global}}$ \\
\midrule
Solo CSP            & $\times$   & \checkmark & $\times$   & $\times$   & 1.000 \\
ML + CSP            & \checkmark & \checkmark & $\times$   & $\times$   & 1.000 \\
ML + CSP + Bayes    & \checkmark & \checkmark & \checkmark & $\times$   & 0.281 \\
Sistema completo    & \checkmark & \checkmark & \checkmark & \checkmark & 0.219 \\
\bottomrule
\multicolumn{6}{l}{\small\textit{Nota}: La reducci\'on al agregar Bayes refleja el riesgo
  real del horario ($P\!=\!0.719$); el MDP pondera adem\'as la calidad de la pol\'itica.}
\end{tabular}
\end{table}

%% ============================================================================
\chapter{Conclusiones}
\label{ch:conclusiones}

\section{Conclusiones}

\begin{enumerate}[leftmargin=2em,itemsep=6pt]

  \item \textbf{Integraci\'on coherente de m\'ultiples t\'ecnicas de IA.}
    El sistema h\'ibrido demuestra que ML, CSP, redes bayesianas y MDP no son
    herramientas aisladas: cada una resuelve una capa espec\'ifica del problema
    y sus salidas se encadenan de forma l\'ogica en el pipeline de decisi\'on.

  \item \textbf{Desempe\~no comprobado del m\'odulo ML.}
    El mejor modelo, Random Forest, alcanz\'o F1\,=\,0.934 y AUC\,=\,0.981
    sobre el dataset OULAD real ($n\!=\!11\,999$), superando a todos los baselines.
    El umbral de decisi\'on \'optimo fue $\tau=0.23$, obtenido minimizando el coste
    diferencial entre falsos positivos y falsos negativos.

  \item \textbf{Efectividad del m\'odulo CSP.}
    El solucionador BT+FC+MRV+LCV resolvi\'o la instancia de 10~secciones
    con \textbf{0 retrocesos} y $\eta_H=1.0$, generando un horario factible en
    tiempo despreciable.  La red bayesiana revel\'o un riesgo operativo elevado
    ($P(\text{Replanif})=0.719$), lo que ilustra el valor de incorporar el
    m\'odulo~3: el CSP en solitario no detecta esta vulnerabilidad.

  \item \textbf{Aplicabilidad real.}
    El problema de planificaci\'on acad\'emica es un caso de uso real en cualquier
    universidad. Los m\'odulos propuestos son implementables con bibliotecas
    abiertas de Python y datos hist\'oricos de inscripci\'on.

  \item \textbf{Cobertura de los temas del curso.}
    El proyecto integra expl\'icitamente: agentes inteligentes, b\'usqueda
    heur\'istica, CSP/factor graphs, beam search, redes bayesianas,
    aprendizaje supervisado y MDP, demostrando una comprensi\'on transversal
    del curso de IA.

\end{enumerate}

\section{Limitaciones}

\begin{itemize}[leftmargin=2em,itemsep=4pt]
  \item La red bayesiana requiere estimaci\'on manual de las tablas de probabilidad
    condicional en ausencia de datos hist\'oricos suficientes.
  \item El MDP simplificado asume un espacio de estados finito y discreto; la
    realidad puede ser m\'as compleja.
  \item La escalabilidad del backtracking disminuye para instancias con $n > 150$
    secciones sin optimizaciones adicionales.
\end{itemize}

%% ===========================================================================

\nocite{*}
\printbibliography[heading=bibintoc,title={Bibliografía}]

\end{document}